% GENERATED FILE. Edit canonical lessons, then run npm run generate:books.
% canonical-source-hash: fnv1a64:d8d03a6dcc2343e2
% canonical-lessons: FA-C238-tord, FA-C238-abdar, FA-C238-konjkav, FA-C238-sakhtgir, FA-C238-zahmatkesh

\chapter{Crisp, Juicy, Curious, Strict, Hard-working}
\label{ch:fa-crisp-juicy-curious-strict-hard-working}
\hlchaptermodality{238}

\begin{chapteropening}
\textbf{By the end of this chapter:} \emph{I can say crisp, juicy, curious, strict and hard-working.}

Complete the last lesson of chapter 238: I can say crisp, juicy, curious, strict and hard-working.
\end{chapteropening}

\section[\texorpdfstring{\fa{ترد} (tord)}{ترد (tord)}]{\fa{ترد} (tord) — crisp, crunchy}

\label{lesson:FA-C238-tord}



\begin{quote}
Before the new one: say the Persian for alive, then the Persian for dead.
\end{quote}

\subsection*{You'll want to know: \fa{ترد}}
\textbf{\fa{ترد}} — \emph{tord} — \textquotedblleft{}crisp, crunchy\textquotedblright{}.

\begin{cousinweb}
\emph{tord} is a native Persian word, but no source used here gives its older form, and its origin is not settled.
\end{cousinweb}

\begin{grammarlens}[title={What you've built}]
One more word for this chapter.
\end{grammarlens}

\subsection*{Your turn}
\begin{itemize}
\raggedright
  \item \emph{Say it:} \emph{tord}
  \item \emph{Say it:} \emph{tord}, once more
  \item \emph{Say it:} say \emph{tord}
  \item \emph{Recall:} say the Persian for alive, then the Persian for dead, then say \emph{tord} again
\end{itemize}

\begin{tcolorbox}[breakable,colback=teal!4,colframe=teal!35!black,title={Before you move on}]
What does \fa{ترد} mean? (\textquotedblleft{}Crisp, crunchy\textquotedblright{}.) Say it once more.
\end{tcolorbox}
\section[\texorpdfstring{\fa{آبدار} (âbdâr)}{آبدار (âbdâr)}]{\fa{آبدار} (âbdâr) — juicy}

\label{lesson:FA-C238-abdar}



\begin{quote}
Before the new one: say the Persian for dead, then the Persian for crisp.
\end{quote}

\subsection*{You'll want to know: \fa{آبدار}}
\textbf{\fa{آبدار}} — \emph{âbdâr} — \textquotedblleft{}juicy\textquotedblright{}.

\begin{cousinweb}
\emph{âbdâr} is \emph{âb}, \textquotedblleft{}water, juice\textquotedblright{}, with \emph{-dâr}, \textquotedblleft{}holding\textquotedblright{}: juice-holding. Both are native; \emph{âb} continues Middle Persian \emph{āb} (EtymDB).
\end{cousinweb}

\begin{grammarlens}[title={What you've built}]
One more word for this chapter.
\end{grammarlens}

\subsection*{Your turn}
\begin{itemize}
\raggedright
  \item \emph{Say it:} \emph{âbdâr}
  \item \emph{Say it:} \emph{âbdâr}, once more
  \item \emph{Say it:} say \emph{âbdâr}
  \item \emph{Recall:} say the Persian for dead, then the Persian for crisp, then say \emph{âbdâr} again
\end{itemize}

\begin{tcolorbox}[breakable,colback=teal!4,colframe=teal!35!black,title={Before you move on}]
What does \fa{آبدار} mean? (\textquotedblleft{}Juicy\textquotedblright{}.) Say it once more.
\end{tcolorbox}
\section[\texorpdfstring{\fa{کنجکاو} (konjkâv)}{کنجکاو (konjkâv)}]{\fa{کنجکاو} (konjkâv) — curious}

\label{lesson:FA-C238-konjkav}



\begin{quote}
Before the new one: say the Persian for crisp, then the Persian for juicy.
\end{quote}

\subsection*{You'll want to know: \fa{کنجکاو}}
\textbf{\fa{کنجکاو}} — \emph{konjkâv} — \textquotedblleft{}curious\textquotedblright{}.

\begin{cousinweb}
\emph{konjkâv} is built from \emph{konj}, \textquotedblleft{}corner\textquotedblright{}, and \emph{-kâv}, \textquotedblleft{}digging\textquotedblright{} (from \emph{kâvidan}, \textquotedblleft{}to dig, to search\textquotedblright{}): one who searches the corners. Both are native Persian words, inherited rather than borrowed.
\end{cousinweb}

\begin{grammarlens}[title={What you've built}]
One more word for this chapter.
\end{grammarlens}

\subsection*{Your turn}
\begin{itemize}
\raggedright
  \item \emph{Say it:} \emph{konjkâv}
  \item \emph{Say it:} \emph{konjkâv}, once more
  \item \emph{Say it:} say \emph{konjkâv}
  \item \emph{Recall:} say the Persian for crisp, then the Persian for juicy, then say \emph{konjkâv} again
\end{itemize}

\begin{tcolorbox}[breakable,colback=teal!4,colframe=teal!35!black,title={Before you move on}]
What does \fa{کنجکاو} mean? (\textquotedblleft{}Curious\textquotedblright{}.) Say it once more.
\end{tcolorbox}
\section[\texorpdfstring{\fa{سختگیر} (sakhtgir)}{سختگیر (sakhtgir)}]{\fa{سختگیر} (sakhtgir) — strict}

\label{lesson:FA-C238-sakhtgir}



\begin{quote}
Before the new one: say the Persian for juicy, then the Persian for curious.
\end{quote}

\subsection*{You'll want to know: \fa{سختگیر}}
\textbf{\fa{سختگیر}} — \emph{sakhtgir} — \textquotedblleft{}strict\textquotedblright{}.

\begin{cousinweb}
\emph{sakhtgir} is \emph{sakht}, \textquotedblleft{}hard\textquotedblright{}, with \emph{-gir}, \textquotedblleft{}taking, gripping\textquotedblright{} (from \emph{gereftan}): one who grips hard. \emph{sakht} is inherited, from Middle Persian \emph{saxt} (EtymDB).
\end{cousinweb}

\begin{grammarlens}[title={What you've built}]
One more word for this chapter.
\end{grammarlens}

\subsection*{Your turn}
\begin{itemize}
\raggedright
  \item \emph{Say it:} \emph{sakhtgir}
  \item \emph{Say it:} \emph{sakhtgir}, once more
  \item \emph{Say it:} say \emph{sakhtgir}
  \item \emph{Recall:} say the Persian for juicy, then the Persian for curious, then say \emph{sakhtgir} again
\end{itemize}

\begin{tcolorbox}[breakable,colback=teal!4,colframe=teal!35!black,title={Before you move on}]
What does \fa{سختگیر} mean? (\textquotedblleft{}Strict\textquotedblright{}.) Say it once more.
\end{tcolorbox}
\section[\texorpdfstring{\fa{زحمتکش} (zahmatkesh)}{زحمتکش (zahmatkesh)}]{\fa{زحمتکش} (zahmatkesh) — hard-working}

\label{lesson:FA-C238-zahmatkesh}



\begin{quote}
Before the new one: say the Persian for curious, then the Persian for strict.
\end{quote}

\subsection*{You'll want to know: \fa{زحمتکش}}
\textbf{\fa{زحمتکش}} — \emph{zahmatkesh} — \textquotedblleft{}hard-working\textquotedblright{}.

\begin{cousinweb}
\emph{zahmatkesh} is \emph{zahmat}, \textquotedblleft{}trouble, toil\textquotedblright{}, with \emph{-kesh}, \textquotedblleft{}bearing\textquotedblright{} (from \emph{keshidan}): one who bears toil. \emph{zahmat} is borrowed from Arabic \emph{zaḥma}, \textquotedblleft{}crowding, pressure\textquotedblright{} (root \emph{z-ḥ-m}, Lane); \emph{kesh} is native.
\end{cousinweb}

\begin{grammarlens}[title={What you've built}]
One more word for this chapter.
\end{grammarlens}

\subsection*{Your turn}
\begin{itemize}
\raggedright
  \item \emph{Say it:} \emph{zahmatkesh}
  \item \emph{Say it:} \emph{zahmatkesh}, once more
  \item \emph{Say it:} say \emph{zahmatkesh}
  \item \emph{Recall:} say the Persian for curious, then the Persian for strict, then say \emph{zahmatkesh} again
\end{itemize}

\begin{tcolorbox}[breakable,colback=teal!4,colframe=teal!35!black,title={Before you move on}]
What does \fa{زحمتکش} mean? (\textquotedblleft{}Hard-working\textquotedblright{}.) Say it once more.
\end{tcolorbox}
